% problem_id: S001-lemma-compact
% source_id: S001 (Stacks Project)
% source_file: dga.tex
% statement_locator: dga.tex lines 6373--6385
% proof_locator: dga.tex lines 6387--6482
% env: lemma
% pairing: tight (verified)
% status: complete (statement + author proof verbatim + reason + locators)
%
\begin{lemma}
\label{lemma-compact}
Let $(A, \text{d})$ be a differential graded algebra. Assume that $A^n = 0$
for $|n| \gg 0$. Let $E$ be an object of $D(A, \text{d})$.
The following are equivalent
\begin{enumerate}
\item $E$ is a compact object, and
\item $E$ can be represented by a differential graded $A$-module $P$
which is finite projective as a graded $A$-module and satisfies
$\Hom_{K(A, \text{d})}(P, M) = \Hom_{D(A, \text{d})}(P, M)$
for every differential graded $A$-module $M$.
\end{enumerate}
\end{lemma}

\begin{proof}
Let $\mathcal{D} \subset K(A, \text{d})$ be the triangulated subcategory
discussed in Remark \ref{remark-source-graded-projective}.
Let $P$ be an object of $\mathcal{D}$ which is finite projective
as a graded $A$-module. Then $P$ represents a compact object of
$D(A, \text{d})$ by Remark \ref{remark-graded-projective-is-compact}.

\medskip\noindent
To prove the converse, let $E$ be a compact object of $D(A, \text{d})$.
Fix $a \leq b$ as in Lemma \ref{lemma-compact-implies-bounded}.
After decreasing $a$ and increasing $b$ if necessary, we may also
assume that $H^i(E) = 0$ for $i \not \in [a, b]$ (this follows
from Proposition \ref{proposition-compact} and our assumption on $A$).
Moreover, fix an integer $c > 0$ such that $A^n = 0$ if $|n| \geq c$.

\medskip\noindent
By Proposition \ref{proposition-compact} we see that $E$ is a direct
summand, in $D(A, \text{d})$, of a differential graded $A$-module $P$
which has a finite filtration $F_\bullet$ by differential
graded submodules such that $F_iP/F_{i - 1}P$ are finite direct sums
of shifts of $A$. In particular, $P$ has property (P) and we have
$\Hom_{D(A, \text{d})}(P, M) = \Hom_{K(A, \text{d})}(P, M)$ for any
differential graded module $M$ by Lemma \ref{lemma-hom-derived}.
In other words, $P$ is an object of the triangulated
subcategory $\mathcal{D} \subset K(A, \text{d})$ discussed in
Remark \ref{remark-source-graded-projective}.
Note that $P$ is finite free as a graded $A$-module.

\medskip\noindent
Choose $n > 0$ such that $b + 4c - n < a$.
Represent the projector onto $E$ by an endomorphism $\varphi : P \to P$ of
differential graded $A$-modules. Consider the distinguished triangle
$$
P \xrightarrow{1 - \varphi} P \to C \to P[1]
$$
in $K(A, \text{d})$ where $C$ is the cone of the first arrow. Then
$C$ is an object of $\mathcal{D}$,
we have $C \cong E \oplus E[1]$ in $D(A, \text{d})$, and
$C$ is a finite graded free $A$-module.
Next, consider a distinguished triangle
$$
C[1] \to C \to C' \to C[2]
$$
in $K(A, \text{d})$ where $C'$ is the cone on a morphism $C[1] \to C$
representing the composition
$$
C[1] \cong E[1] \oplus E[2] \to E[1] \to E \oplus E[1] \cong C
$$
in $D(A, \text{d})$. Then we see that $C'$ represents $E \oplus E[2]$.
Continuing in this manner we see that we can find a differential
graded $A$-module $P$ which is an object of $\mathcal{D}$,
is a finite free as a graded $A$-module, and represents $E \oplus E[n]$.

\medskip\noindent
Choose a basis $x_i$, $i \in I$ of homogeneous elements for $P$ as an
$A$-module. Let $d_i = \deg(x_i)$.
Let $P_1$ be the $A$-submodule of $P$ generated by $x_i$ and
$\text{d}(x_i)$ for $d_i \leq a - c - 1$.
Let $P_2$ be the $A$-submodule of $P$ generated by $x_i$ and
$\text{d}(x_i)$ for $d_i \geq b - n + c$.
We observe
\begin{enumerate}
\item $P_1$ and $P_2$ are differential graded submodules of $P$,
\item $P_1^t = 0$ for $t \geq a$,
\item $P_1^t = P^t$ for $t \leq a - 2c$,
\item $P_2^t = 0$ for $t \leq b - n$,
\item $P_2^t = P^t$ for $t \geq b - n + 2c$.
\end{enumerate}
As $b - n + 2c \geq a - 2c$ by our choice of $n$
we obtain a short exact sequence of differential graded $A$-modules
$$
0 \to P_1 \cap P_2 \to P_1 \oplus P_2 \xrightarrow{\pi} P \to 0
$$
Since $P$ is projective as a graded $A$-module this is an admissible
short exact sequence (Lemma \ref{lemma-target-graded-projective}).
Hence we obtain a boundary map
$\delta : P \to (P_1 \cap P_2)[1]$ in $K(A, \text{d})$, see
Lemma \ref{lemma-admissible-ses}.
Since $P = E \oplus E[n]$ and since $P_1 \cap P_2$ lives in
degrees $(b - n, a)$ we find that
$\Hom_{D(A, \text{d})}(E \oplus E[n], (P_1 \cap P_2)[1])$ is
zero. Therefore $\delta = 0$ as a morphism in $K(A, \text{d})$
as $P$ is an object of $\mathcal{D}$.
By Derived Categories, Lemma \ref{derived-lemma-split}
we can find a map $s : P \to P_1 \oplus P_2$ such that
$\pi \circ s = \text{id}_P + \text{d}h + h\text{d}$ for some $h : P \to P$
of degree $-1$. Since $P_1 \oplus P_2 \to P$ is surjective and since $P$
is projective as a graded $A$-module we can choose a homogeneous
lift $\tilde h : P \to P_1 \oplus P_2$ of $h$. Then we change
$s$ into $s + \text{d} \tilde h + \tilde h \text{d}$ to get
$\pi \circ s = \text{id}_P$. This means we obtain a direct
sum decomposition $P = s^{-1}(P_1) \oplus s^{-1}(P_2)$.
Since $s^{-1}(P_2)$ is equal to $P$ in degrees $\geq b - n + 2c$
we see that $s^{-1}(P_2) \to P \to E$ is a quasi-isomorphism,
i.e., an isomorphism in $D(A, \text{d})$. This finishes the proof.
\end{proof}
